\begin{abstract}
本文针对多波束测深测线布设问题，建立了覆盖宽度、视坡度和测线优化模型。问题1中，基于二维剖面几何关系推导覆盖宽度与重叠率公式；当水深为70m、坡度为$1.5^\circ$时，覆盖宽度为243.07m，平坦近似宽度为242.49m，相对差异为0.24\%。问题2中，引入测线方向与坡面法向水平投影夹角$\beta$，建立视坡度模型 $\alpha'=\arctan(\tan\alpha\cos\beta)$；验证得到$\beta=30^\circ$时视坡度为8.6822^\circ$，$\beta=0^\circ$时为10.0^\circ$，$\beta=90^\circ$时为0.0^\circ$；result2 表中覆盖宽度范围为63.05m--770.33m。

问题3中，将恒定坡度矩形海域的布线转化为一维变间距优化问题，以全覆盖和重叠率约束为条件，采用贪心迭代确定相邻测线位置。结果得到34条测线，总长度为125936.0m，相邻重叠率范围为10.5\%--10.5\%，满足10\%--20\%约束。问题4中，针对真实水深网格，建立基于动态覆盖宽度的自适应贪心规划模型，以候选测线对未覆盖网格的覆盖面积增量为收益函数。最终布设89条测线，总长度为558197.66m，漏测率为1.8771\%，超重叠率区域长度为142834.27m，占总测线长度25.5885\%。灵敏度分析表明，坡度增加10\%使总长度增加38.24\%，开角增加20\%使总长度减少44.12\%，目标重叠率增加20\%使总长度增加2.94\%。

\textbf{关键词}：多波束测深；覆盖宽度模型；视坡度；测线优化；自适应贪心算法；重叠率；灵敏度分析
\end{abstract}